\section{Conclusion}

One premise, taken seriously, forces almost the whole architecture of physics. There cannot be nothing, so there is one
distinction, three-valued because a distinction has three parts. Combining distinctions without loss forces the
octonions, and the octonions force four-dimensional spacetime, the 24-cell, the hyperbolic mesh, and a single
reversible law. Off that one structure come spinors and one generation of quarks and leptons with exact charges, the
three forces and the Higgs, emergent relativity and gravity and the quantum, holography and black-hole thermodynamics,
three spatial dimensions and a constant dark energy, bound matter and a dark-matter candidate, and a universal computer
with a solved address for every cell. Each step is the unique answer to a short requirement, and each is checked by a
runnable experiment with a control where the answer should come out no.

The absolute masses and couplings are free, exactly the parameters the Standard Model already leaves free, each now
carrying a geometric origin, so the model does not claim to derive every number, it claims that the architecture, the
part that always looked arbitrary, is forced. One deep question stays open and is named loudly rather than buried,
whether the forced threefold structure is the three generations of matter we observe, and a fundamental reversible
gravity, a handful of computations, and a few statistical estimators remain frontiers. Where the base gives the physics,
the paper shows it and names the experiment. Where it only reproduces what is known, it says so. Where it fails or
stays open, it says that too.

This is exploratory, feasibility-stage work, offered as a direction worth pursuing, not a finished theory. Its claim is
modest in form and large in reach: that a single discrete, deterministic, reversible rule, on a single forced mesh,
made of one substance and nothing else, reproduces a long and specific list of the physics we measure, and is itself a
universal computer. Whether the universe began or is eternal is left open, and whether the substrate truly feels is the
one posit, set aside and unused in everything above. What remains, when the posit is set down, is a discrete model of the
universe that runs, that derives what it can, that is plain about what it cannot, and that could be shown wrong, which is
all a theory can ask to be. The full derivation, from nothing to the cell, is laid out once more in the appendix, slowly,
for anyone who wants to walk it from the ground up.
